\documentclass[11pt,a4paper]{article}
\usepackage[margin=25mm]{geometry}
\usepackage{amsmath,amssymb,amsthm}
\newtheorem{proposition}{Proposition}
\title{A qualification to an orthogonal rank bound}
\author{Human Compression project: F001 source audit}\date{}
\begin{document}\maketitle
Brooksbank, \emph{On the ranks of string C-group representations for
symplectic and orthogonal groups}, arXiv:1912.05512v1, physical page~8,
Lemma~4.1, asserts that a sequence of at least six involutions in a
subgroup of a five-dimensional orthogonal group satisfying the string
commutations must fail the intersection property. Its stated scope
includes arbitrary subgroups of the orthogonal group. The following
standard permutation-module construction is a counterexample within
that scope. It does not disprove the principal theorem concerning
$\operatorname{PSp}_4(q)$.

\begin{proposition}
There is a nondegenerate quadratic space $V$ of dimension five over
$\mathbb F_7$ and a subgroup $H\le O(V)$ having a connected string
C-marking of rank six, with every adjacent product of order three.
One may moreover take $H\le SO(V)$.
\end{proposition}
\begin{proof}
On $M=\mathbb F_7^7$ use the coordinate dot product and quadratic form
$Q(x)=\sum_{i=1}^7x_i^2$. Put
\[
 W=\{x:\textstyle\sum_i x_i=0\},\qquad
 T=\langle(1,1,1,1,1,1,1)\rangle.
\]
The dot product on $M$ is nondegenerate, $W=T^\perp$, and
$T\subset W$. Hence the radical of its restriction to $W$ is exactly
$T$. Since $Q$ vanishes on $T$ and $T$ is orthogonal to $W$, the
form descends to a nondegenerate quadratic form on $V=W/T$, of
dimension five. Coordinate permutations preserve $W,T,Q$ and give
an orthogonal representation $\pi:S_7\longrightarrow O(V)$.

This representation is faithful. If $\pi(g)=1$, then for every
$i\ne j$ the vector
$e_{g(i)}-e_{g(j)}-e_i+e_j$ belongs to $T$. It has at most four
nonzero coordinates; a nonzero vector of $T$ has all seven coordinates
nonzero. Thus it is zero. Since the characteristic is not two, this
forces $g(i)=i$ and $g(j)=j$. Applying this to all pairs gives $g=1$.
The same support argument shows $-1_V\notin\pi(S_7)$: otherwise
$g$ would interchange $i,j$ for every distinct pair, which is impossible.

Let $s_i=(i\ i+1)$ for $1\le i\le6$. The tuple
$(\pi(s_1),\ldots,\pi(s_6))$ generates $\pi(S_7)$, consists of
involutions and has the required string commutations and adjacent
orders three. For completeness, its intersection property can be
checked before applying the faithful representation. For any
$I\subseteq\{1,\ldots,6\}$, the subgroup generated by $s_i$, $i\in I$,
is the full product of symmetric groups on the interval components
of the path with edge set $I$. Two such subgroups intersect in the
product of symmetric groups on the common refinement of their
interval partitions. That refinement consists exactly of the
components of the path with edge set $I\cap J$. Therefore
\[
 \langle s_i:i\in I\rangle\cap\langle s_j:j\in J\rangle
       =\langle s_k:k\in I\cap J\rangle.
\]
Faithfulness transfers this equality to the displayed marking in $O(V)$.

Alternatively use $\widetilde\pi(g)=\operatorname{sgn}(g)\pi(g)$.
This is again a representation; its kernel is trivial because $\pi$
is faithful and its image does not contain $-1_V$. Each $\pi(s_i)$
acts by $-1$ on the image of $e_i-e_{i+1}$ and fixes its orthogonal
hyperplane, so has determinant $-1$. In dimension five its negative
has determinant one. The signed representation is consequently in
$SO(V)$ and transfers exactly the same string C-marking.
\end{proof}

\paragraph{Coordinates and the precise failed inference.}
A basis of $V$ is given by the classes of $e_i-e_7$, $1\le i\le5$.
Its dot-product Gram matrix is $B=I_5+J_5$, of determinant
$6\in\mathbb F_7^\times$. The supplied six literal matrices are
computed by permuting these vectors and reducing modulo $T$.
The accompanying exact program checks their group order $5040$,
all isometry and string identities, and all $64^2$ subset
intersection pairs, for both the ordinary and signed representations.
These computations corroborate the elementary proof above.

The source proof passes from a nonzero intersection of linear supports
to a nontrivial intersection of generated groups. In the ordinary
representation, the first and last triples each have a three-dimensional
support, and their supports intersect in one dimension. Their groups
nevertheless intersect trivially, by the displayed intersection property.
Thus that implication cannot justify the stated orthogonal bound.

\paragraph{Scope and attribution.}
The fully deleted permutation module and adjacent-transposition
marking are standard constructions; no priority for either is claimed.
This note concerns the exact bound source edition and one auxiliary
lemma. It supplies neither a counterexample to the $\operatorname{PSp}_4(q)$
rank spectrum nor a replacement proof of that spectrum. In particular,
membership of the signed image in $SO(V)$ does not assert membership
in $\Omega(V)$. The latter source theorem, its uniform parameter
existence argument and its complete recovery remain pending. There is
no open-conjecture credit, compression measurement or Lean claim here.
\end{document}
